% Incremental evidence audit; not a manuscript novelty claim.
\section{2026年10月11日增量核查：剩余问题与基线诊断}
本次沿用已有文献库中的20篇正式会议或期刊方法论文，不新增重复BibTeX条目。综述用于方法分类，剩余问题优先由2025--2026年原文支持。作者明确局限、方法适用范围与本研究推论分别记录；本轮不提前确定创新方法。
\subsection{已经改进与仍需检验的问题}
\begin{enumerate}
\item \textbf{强条件下的联合质量。} 任意关节控制、噪声优化、检索和物理精化已经有方法。RG-DNO仍报告不自然动作与文本偏差，RobotMDM没有硬可行性保证，PARC的混合消融显示接触符合与连续性的权衡。不能把加入条件、物理信息或闭环本身称作创新。\citep{rgdno,robotmdm,parc}
\item \textbf{修正结果的局部验收与适用范围。} Morph已有MPJPE门控及语义、跨任务分析；PARC、CLAIMS与LIMMT已有质量评价或筛选。剩余候选问题是残余局部缺陷、不适用输入与修复失败的辨识，需要逐片段实验验证，不能声称首次数据筛选。\citep{morph,parc,claims,limmt}
\item \textbf{执行反馈和难度的客观校准。} CLAIMS已有等数据预算、多跟踪器验证，LIMMT已验证高质量子集有效。CLAIMS明确指出人工变量库缺少客观校准与全面覆盖；LIMMT规则和选择比例仍需更多语料检验。执行失败同时受参考、场景和控制器影响，这是待验证的归因问题，不能将失败一概解释为错误数据。\citep{claims,limmt}
\end{enumerate}
\subsection{固定基线与测量边界}
使用发布版CLoSD的生成模型与跟踪策略，在已有Isaac Gym环境中运行三个随机种子（11、12、13），每个64环境，全部保留第一次尝试和失败前缀，不更新模型、策略或奖励。记录策略当前观测中应用动作前的生成参考，以及同一步物理执行后的状态。初始30帧仅从质量统计排除，终止统计保留全部帧。\citep{closd}
首轮v2错误读取了继承奖励函数使用的独立数据缓存，其跟踪分析已作废并保留记录。v3改用实际观测参考，通过独立调用顺序、坐标、关节顺序和时间变化检查。这是测量修正，不是控制算法改进。
本轮为SMPL物理角色的文本到动作诊断，不是目标人形机器人执行器认证，不含显式空间目标，也没有验证下游再训练收益。足部刚体中心速度只能作接触代理，旋转可能使其移动，不能直接当作表面滑移真值。
\subsection{实际结果与假设判定}
\begin{center}\begin{tabular}{lrrrr}\hline
种子 & 首次尝试 & 未提前终止 & 误差门控通过 & 联合代理门控通过\\\hline
11 & 64 & 59 & 54 & 29\\
12 & 64 & 57 & 52 & 29\\
13 & 64 & 58 & 53 & 27\\
合计 & 192 & 174 & 159 & 85\\
\hline\end{tabular}\end{center}
预设误差门槛为平均g-MPJPE不超过100毫米、帧误差p95不超过200毫米；联合门控增加持续净接触力超过10牛顿时，足部中心速度超过0.5米/秒的比例不超过5\%。阈值为本次诊断选择，不是通用合格标准。
\begin{itemize}
\item \textbf{H1：有限支持。} 174条未提前终止的片段中15条未通过误差门控。没有提前终止不能替代跟踪质量检查，但尚不能将误差归因于生成器。
\item \textbf{H2：跟踪恶化假设未得到支持。} 参考边界速度比内部高约2.202米/秒，实际运动差约0.058米/秒，边界跟踪误差反而平均低34.64毫米。已执行状态参与后续规划与坐标对齐；参考更新不等同于身体运动不连续。
\item \textbf{接触代理尚不能认证物理缺陷。} 27组阈值敏感性设置下联合通过63--109条，主设置为85条。应结合表面接触几何复核，不能把门控拒绝的样本全部解释为脚滑。
\end{itemize}
原地跳高样本（种子12/环境3）的平均误差120.9毫米、p95为341.2毫米；舞蹈样本（11/61）分别为123.6与342.5毫米。含台阶和扶栏的提示缺少对应场景支撑，应与生成缺陷分开。本轮未支持三项缺口全部成立，也未证明合成数据提高模仿学习。
\subsection{下一步的证据要求}
先在相同输入和条件上比较普通生成与通用物理精化，检查关键事件、跟踪误差和接触几何；排除缺少场景支撑的请求。只有确认具体失败模式后，才确定新的生成或修正方法。进一步在同一片段上使用第二个固定跟踪器分离控制器影响，最后以等数据量、等训练预算、多种子和独立测试条件验证训练价值。当前结论限定于仿真。
